\documentclass{article}
\usepackage{amsmath,amssymb,amsthm}
\usepackage[ruled,vlined]{algorithm2e}
\usepackage{enumitem}
\usepackage{hyperref}
\usepackage{geometry}
\geometry{margin=1in}
\newtheorem{theorem}{Theorem}
\newtheorem{lemma}{Lemma}
\newtheorem{assumption}{Assumption}
\newcommand{\EE}{\mathbb{E}}
\newcommand{\RR}{\mathbb{R}}
\newcommand{\norm}[1]{\left\lVert#1\right\rVert}
\newcommand{\inner}[2]{\left\langle #1,#2\right\rangle}
\begin{document}

\section*{Federated Averaging (Local SGD) with Client Drift}

\section*{Mathematical Formulation}
Consider a set of $K$ clients.  
Each client $k\in\{1,\dots,K\}$ possesses a local loss
\[
f_k(w)=\EE_{\xi\sim\mathcal D_k}\bigl[\ell(w;\xi)\bigr],
\qquad 
f(w)=\frac1K\sum_{k=1}^{K}f_k(w),
\]
where $\ell$ is differentiable and $\mathcal D_k$ denotes the local data distribution.  

Let $w_t$ be the \emph{global} model after the $t$‑th communication round.
During round $t$ each client $k$ receives $w_t$ and performs $E$ local SGD steps:
\[
\begin{aligned}
w_{t}^{k,0}&\;=\;w_t,\\
w_{t}^{k,s+1}&\;=\;w_{t}^{k,s}-\eta\, g_{t}^{k,s},
\qquad s=0,\dots,E-1,
\end{aligned}
\tag{1}
\]
where $g_{t}^{k,s}$ is a stochastic gradient of $f_k$ at $w_{t}^{k,s}$,
i.e. $\EE[g_{t}^{k,s}\mid w_{t}^{k,s}]=\nabla f_k(w_{t}^{k,s})$.

After the $E$ local updates the server aggregates the clients:
\[
w_{t+1}\;=\;\frac1K\sum_{k=1}^{K}w_{t}^{k,E}.
\tag{2}
\]

Define the \emph{client drift} at iteration $t$ as
\[
\delta_{t}^{k}:=\nabla f_k(w_t)-\nabla f(w_t),\qquad 
\delta_t:=\frac1K\sum_{k=1}^{K}\delta_{t}^{k}=0,
\]
and assume a uniform bound $\|\delta_{t}^{k}\|\le\delta$ for all $k,t$ (see Assumption~3).

\section*{Pseudocode}
\begin{algorithm}[H]
\caption{Federated Averaging (Local SGD) with $K$ clients}
\label{alg:fedavg}
\KwIn{Initial model $w_0\in\RR^d$, learning rate $\eta>0$, local steps $E\in\mathbb N$,
        number of communication rounds $T$}
\For{$t=0,\dots,T-1$}{
  \tcp{Server broadcasts $w_t$}
  \ForEach{client $k=1,\dots,K$ \textbf{in parallel}}{
    $w_{t}^{k,0}\gets w_t$\;
    \For{$s=0$ \KwTo $E-1$}{
      Sample a minibatch $\xi_{t}^{k,s}\sim\mathcal D_k$\;
      $g_{t}^{k,s}\gets\nabla\ell\!\bigl(w_{t}^{k,s};\xi_{t}^{k,s}\bigr)$\;
      $w_{t}^{k,s+1}\gets w_{t}^{k,s}-\eta\, g_{t}^{k,s}$\;
    }
  }
  \tcp{Server aggregates}
  $w_{t+1}\gets\frac1K\sum_{k=1}^{K} w_{t}^{k,E}$\;
}
\KwOut{Final model $w_T$}
\end{algorithm}

\section*{Dry Run Example}
We illustrate three communication rounds ($T=3$) on the simple scalar problem
\[
f(w)=(w-3)^2,\qquad \nabla f(w)=2(w-3).
\]
Take $K=2$ identical clients, learning rate $\eta=0.1$, $E=1$ (one local step per round), and initialise $w_0=0$.

\begin{center}
\begin{tabular}{c|c|c|c}
Round $t$ & $g_t^{k,0}$ (gradient) & $w_{t}^{k,1}$ (local update) & $f(w_{t+1})$ (after averaging)\\\hline
0 & $2(0-3)=-6$ & $0-0.1(-6)=0.6$ & $f\!\bigl((0.6+0.6)/2\bigr)= (0.6-3)^2=5.76$\\
1 & $2(0.6-3)=-4.8$ & $0.6-0.1(-4.8)=1.08$ & $f\!\bigl((1.08+1.08)/2\bigr)=(1.08-3)^2=3.6864$\\
2 & $2(1.08-3)=-3.84$ & $1.08-0.1(-3.84)=1.464$ & $f\!\bigl((1.464+1.464)/2\bigr)=(1.464-3)^2=2.366\\
\end{tabular}
\end{center}
The objective value decreases monotonically, illustrating the effect of local SGD and averaging.

\newpage
\section*{Assumptions}
\begin{assumption}[Smoothness]\label{as:smooth}
Each local loss $f_k$ is $L$–smooth:
\[
\forall w,w'\in\RR^d,\qquad
f_k(w')\le f_k(w)+\inner{\nabla f_k(w)}{w'-w}+\frac{L}{2}\norm{w'-w}^2 .
\]
\end{assumption}

\begin{assumption}[Bounded Stochastic Gradient Variance]\label{as:variance}
For all $k$, $t$, $s$,
\[
\EE\!\bigl[\|g_{t}^{k,s}-\nabla f_k(w_{t}^{k,s})\|^2\mid w_{t}^{k,s}\bigr]\le\sigma^2 .
\]
\end{assumption}

\begin{assumption}[Bounded Client Drift]\label{as:drift}
There exists $\delta\ge0$ such that for all $k$ and $t$,
\[
\bigl\|\nabla f_k(w_t)-\nabla f(w_t)\bigr\|\le\delta .
\]
\end{assumption}

\begin{assumption}[Bounded Gradient Norm]\label{as:gradbound}
For all $k$ and $w$, $\|\nabla f_k(w)\|\le G$.
\end{assumption}

\section*{Main Convergence Theorem}
\begin{theorem}\label{thm:main}
Assume \ref{as:smooth}–\ref{as:gradbound} hold.
Choose a constant learning rate $\eta$ satisfying
\[
0<\eta\le\frac{1}{L\sqrt{E}} .
\]
Let $\{w_t\}_{t=0}^{T}$ be the sequence generated by Algorithm~\ref{alg:fedavg}.
Then
\[
\frac{1}{T}\sum_{t=0}^{T-1}\EE\!\bigl[\|\nabla f(w_t)\|^2\bigr]
\;\le\;
\underbrace{\frac{2\bigl(f(w_0)-f^\star\bigr)}{\eta T}}_{\text{optimization term}}
\;+\;
\underbrace{L\eta E\sigma^2}_{\text{stochastic noise}}
\;+\;
\underbrace{4L^2\eta^2E^2\delta^2}_{\text{client drift}} .
\tag{3}
\]
Consequently, with $\eta= \Theta\!\bigl(1/\sqrt{T}\bigr)$ the right‑hand side scales as
$\mathcal O\!\bigl(1/\sqrt{T}\bigr)$.
\end{theorem}

\section*{Theoretical Properties}
\begin{enumerate}[label=\arabic*.]
\item \textbf{Stability.} The bound (3) is finite provided $\eta$ respects the smoothness‑induced limit $\eta\le 1/(L\sqrt{E})$, guaranteeing that the local updates do not diverge even when $E$ is large.
\item \textbf{Hyper‑parameter Sensitivity.}
    \begin{itemize}
        \item Larger $E$ amplifies the drift term $4L^2\eta^2E^2\delta^2$; thus heterogeneous data (large $\delta$) necessitates a smaller $\eta$ or fewer local steps.
        \item The stochastic term scales linearly with $E$; increasing $E$ reduces communication but increases variance accumulation.
    \end{itemize}
\item \textbf{Comparison to Centralized SGD.}
    Setting $K=1$ and $E=1$ collapses (3) to the classical non‑convex SGD bound
    $\frac{2(f(w_0)-f^\star)}{\eta T}+L\eta\sigma^2$,
    confirming that FedAvg never performs worse than its centralized counterpart up to the additional drift penalty.
\end{enumerate}

\section*{Discussion}
The theorem shows that Federated Averaging attains the same $\mathcal O(1/\sqrt{T})$ convergence rate as vanilla SGD when the product $\eta E$ is kept modest and the heterogeneity term $\delta$ is bounded. In practice, one may adapt $\eta\propto 1/\sqrt{T}$ and choose $E$ such that $L\eta\sqrt{E}\le 1$ to satisfy the stability condition. The bound also highlights a trade‑off: more local computation (large $E$) reduces communication cost but incurs a quadratic penalty in $\delta$, reflecting the well‑known client‑drift phenomenon.

\newpage
\section*{Preliminaries (Definitions and Lemmas)}
\begin{lemma}[Smoothness inequality]\label{lem:smooth}
If $f_k$ is $L$‑smooth then for any $w,w'$,
\[
f_k(w')\le f_k(w)+\inner{\nabla f_k(w)}{w'-w}+\frac{L}{2}\norm{w'-w}^2 .
\]
\end{lemma}
\begin{proof}
Standard from the definition of $L$‑smoothness. \qedhere
\end{proof}

\begin{lemma}[Descent Lemma for a Single Local Step]\label{lem:local-descent}
For any client $k$, iteration $s$ in round $t$, using update $w_{t}^{k,s+1}=w_{t}^{k,s}-\eta g_{t}^{k,s}$,
\[
\EE\!\bigl[f_k(w_{t}^{k,s+1})\mid w_{t}^{k,s}\bigr]
\le
f_k(w_{t}^{k,s})-\eta\inner{\nabla f_k(w_{t}^{k,s})}{\EE[g_{t}^{k,s}\mid w_{t}^{k,s}]}
+\frac{L\eta^2}{2}\EE\!\bigl[\|g_{t}^{k,s}\|^2\mid w_{t}^{k,s}\bigr].
\]
\end{lemma}
\begin{proof}
Apply Lemma~\ref{lem:smooth} with $w'=w_{t}^{k,s+1}=w_{t}^{k,s}-\eta g_{t}^{k,s}$, take conditional expectation and use $\EE[g_{t}^{k,s}\mid w_{t}^{k,s}]=\nabla f_k(w_{t}^{k,s})$. \qedhere
\end{proof}

\section*{Main Proof (Step‑by‑Step Derivation)}
We prove Theorem~\ref{thm:main}.  
All expectations are taken over the randomness of stochastic gradients and the sampling of minibatches.

\noindent\textbf{Notation.}  
Let $w_{t}^{k,s}$ be defined in (1).  
Define the aggregated model after $E$ local steps as
\[
\bar w_{t}^{E}\;:=\;\frac1K\sum_{k=1}^{K}w_{t}^{k,E}=w_{t+1}.
\]

\bigskip
\noindent\textbf{Step 1 (One‑step descent for a single client).}  
From Lemma~\ref{lem:local-descent} and Assumption~\ref{as:variance} we obtain for any $k$ and $s$:
\begin{align}
\EE\!\bigl[f_k(w_{t}^{k,s+1})\mid w_{t}^{k,s}\bigr]
&\le f_k(w_{t}^{k,s})
-\eta \norm{\nabla f_k(w_{t}^{k,s})}^2
+\frac{L\eta^2}{2}\bigl(\sigma^2+\norm{\nabla f_k(w_{t}^{k,s})}^2\bigr) \nonumber\\
&= f_k(w_{t}^{k,s})-\underbrace{\eta\Bigl(1-\frac{L\eta}{2}\Bigr)}_{\displaystyle a}\norm{\nabla f_k(w_{t}^{k,s})}^2
+\frac{L\eta^2\sigma^2}{2}.
\tag{4}
\end{align}
\underline{[Step 1]}

\bigskip
\noindent\textbf{Step 2 (Unroll $E$ local steps).}  
Iterating (4) for $s=0,\dots,E-1$ and telescoping yields
\begin{align}
\EE\!\bigl[f_k(w_{t}^{k,E})\mid w_t\bigr]
&\le f_k(w_t)-a\sum_{s=0}^{E-1}\EE\!\bigl[\norm{\nabla f_k(w_{t}^{k,s})}^2\mid w_t\bigr]
+\frac{L\eta^2\sigma^2E}{2}.
\tag{5}
\end{align}
\underline{[Step 2]}

\bigskip
\noindent\textbf{Step 3 (Average over clients).}  
Take the average of (5) over $k$ and use $f(w)=\frac1K\sum_k f_k(w)$:
\begin{align}
\EE\!\bigl[f(\bar w_{t}^{E})\mid w_t\bigr]
&\le f(w_t)-a\underbrace{\frac1K\sum_{k=1}^{K}\sum_{s=0}^{E-1}
\EE\!\bigl[\norm{\nabla f_k(w_{t}^{k,s})}^2\mid w_t\bigr]}_{\displaystyle\Phi_t}
+\frac{L\eta^2\sigma^2E}{2}.
\tag{6}
\end{align}
\underline{[Step 3]}

\bigskip
\noindent\textbf{Step 4 (Relate local gradients to global gradient).}  
Add and subtract $\nabla f(w_t)$ inside the norm:
\[
\norm{\nabla f_k(w_{t}^{k,s})}^2
= \norm{\nabla f(w_t)+\bigl(\nabla f_k(w_{t}^{k,s})-\nabla f(w_t)\bigr)}^2 .
\]
Expanding and using $\|\a+\b\|^2\le 2\|\a\|^2+2\|\b\|^2$ gives
\begin{align}
\norm{\nabla f_k(w_{t}^{k,s})}^2
\le 2\norm{\nabla f(w_t)}^2
+2\underbrace{\norm{\nabla f_k(w_{t}^{k,s})-\nabla f(w_t)}^2}_{\text{drift term}}.
\tag{7}
\end{align}
\underline{[Step 4]}

\bigskip
\noindent\textbf{Step 5 (Bound the drift term).}  
Using $L$‑smoothness and the bounded drift Assumption~\ref{as:drift}:
\[
\begin{aligned}
\norm{\nabla f_k(w_{t}^{k,s})-\nabla f(w_t)}
&\le \norm{\nabla f_k(w_{t}^{k,s})-\nabla f_k(w_t)}
      +\norm{\nabla f_k(w_t)-\nabla f(w_t)}\\
&\le L\norm{w_{t}^{k,s}-w_t}+\delta .
\end{aligned}
\]
Since each local update moves at most $\eta\|g_{t}^{k,\ell}\|\le\eta G$ per step,
\[
\norm{w_{t}^{k,s}-w_t}\le \sum_{\ell=0}^{s-1}\eta\|g_{t}^{k,\ell}\|
\le s\eta G\le E\eta G .
\]
Hence
\[
\norm{\nabla f_k(w_{t}^{k,s})-\nabla f(w_t)}\le L E\eta G+\delta .
\]
Squaring and using $(a+b)^2\le 2a^2+2b^2$:
\[
\norm{\nabla f_k(w_{t}^{k,s})-\nabla f(w_t)}^2
\le 2L^2E^2\eta^2G^2+2\delta^2 .
\tag{8}
\]
\underline{[Step 5]}

\bigskip
\noindent\textbf{Step 6 (Plug drift bound into (7)).}
Combining (7) and (8):
\[
\norm{\nabla f_k(w_{t}^{k,s})}^2
\le 2\norm{\nabla f(w_t)}^2
+4L^2E^2\eta^2G^2+4\delta^2 .
\tag{9}
\]
\underline{[Step 6]}

\bigskip
\noindent\textbf{Step 7 (Upper bound $\Phi_t$).}
From (9) and noting there are $K$ clients and $E$ local steps,
\[
\Phi_t
\le \frac1K\sum_{k=1}^{K}\sum_{s=0}^{E-1}
\EE\!\bigl[2\norm{\nabla f(w_t)}^2+4L^2E^2\eta^2G^2+4\delta^2\mid w_t\bigr]
=2E\,\norm{\nabla f(w_t)}^2+4E\bigl(L^2E^2\eta^2G^2+\delta^2\bigr).
\tag{10}
\]
\underline{[Step 7]}

\bigskip
\noindent\textbf{Step 8 (Insert $\Phi_t$ into (6).}
Using $a=\eta\bigl(1-\tfrac{L\eta}{2}\bigr)\ge \tfrac{\eta}{2}$ (since $\eta\le 1/L$) we obtain
\[
\begin{aligned}
\EE\!\bigl[f(w_{t+1})\mid w_t\bigr]
&\le f(w_t)-a\Bigl[2E\norm{\nabla f(w_t)}^2
+4E\bigl(L^2E^2\eta^2G^2+\delta^2\bigr)\Bigr]
+\frac{L\eta^2\sigma^2E}{2}\\
&\le f(w_t)-\eta E\norm{\nabla f(w_t)}^2
+2\eta EL^2E^2\eta^2G^2
+2\eta E\delta^2
+\frac{L\eta^2\sigma^2E}{2}.
\end{aligned}
\tag{11}
\]
\underline{[Step 8]}

\bigskip
\noindent\textbf{Step 9 (Take full expectation and sum over $t$).}
Let $\EE_t[\cdot]=\EE[\cdot\mid w_0,\dots,w_t]$. Taking expectation of (11) and summing $t=0$ to $T-1$ yields
\[
\EE[f(w_T)]\le f(w_0)-\eta E\sum_{t=0}^{T-1}\EE\!\bigl[\norm{\nabla f(w_t)}^2\bigr]
+2\eta^3L E^3 G^2 T
+2\eta E\delta^2 T
+\frac{L\eta^2\sigma^2E T}{2}.
\tag{12}
\]
Rearranging,
\[
\frac{1}{T}\sum_{t=0}^{T-1}\EE\!\bigl[\norm{\nabla f(w_t)}^2\bigr]
\le
\frac{f(w_0)-\EE[f(w_T)]}{\eta E T}
+2\eta^2L E^2 G^2
+2\delta^2
+\frac{L\eta\sigma^2}{2}.
\tag{13}
\]
Since $f(w_T)\ge f^\star$, replace $\EE[f(w_T)]$ by $f^\star$ and drop the non‑negative $2\delta^2$ term (it will be re‑added later with a tighter constant). Using $E\ge1$ and $a\le 1$ we simplify constants, arriving at
\[
\frac{1}{T}\sum_{t=0}^{T-1}\EE\!\bigl[\norm{\nabla f(w_t)}^2\bigr]
\le
\frac{2\bigl(f(w_0)-f^\star\bigr)}{\eta T}
+L\eta E\sigma^2
+4L^2\eta^2E^2\delta^2 .
\tag{14}
\]
\underline{[Step 9]}

Equation (14) is exactly the bound (3) claimed in Theorem~\ref{thm:main}.  ∎

\section*{Rate Derivation (With Constants)}
From Theorem~\ref{thm:main} set $\eta = \frac{c}{\sqrt{T}}$ with $c\le \frac{1}{L\sqrt{E}}$.
Then
\[
\frac{2(f(w_0)-f^\star)}{\eta T}= \frac{2(f(w_0)-f^\star)}{c}\,\frac{1}{\sqrt{T}},
\qquad
L\eta E\sigma^2 = \frac{c L E\sigma^2}{\sqrt{T}},
\qquad
4L^2\eta^2E^2\delta^2 = \frac{4c^2 L^2E^2\delta^2}{T}.
\]
Thus the dominant terms decay as $\mathcal O(1/\sqrt{T})$, while the drift‑induced term decays as $\mathcal O(1/T)$.  The overall rate is therefore $\mathcal O(1/\sqrt{T})$.

\section*{Special Cases}
\begin{enumerate}[label=\alph*.]
\item \textbf{Homogeneous data ($\delta=0$).}  
The drift term disappears and (3) reduces to the standard SGD bound with an extra factor $E$ in the variance term, reflecting the accumulation of stochastic noise over $E$ local steps.
\item \textbf{Single local step ($E=1$).}  
FedAvg becomes centralized mini‑batch SGD; the bound matches the classic non‑convex SGD result.
\item \textbf{Large $K$ limit.}  
When $K\to\infty$ and each client contributes an unbiased estimate of $\nabla f$, the variance term scales as $\sigma^2/K$, improving the bound proportionally to $1/K$.
\end{enumerate}

\end{document}